% Referências do artigo (estilo do modelo IJRER: numeradas na ordem da primeira citação).
% A ordem aqui não importa: o monta.py reordena pelas citações do texto e grava refs_ordenadas.tex.
% Fonte de cada dado: lista de referências do memorial de cálculo e aba Fontes da planilha de precificação.

\bibitem{oliveira2015} A. P. N. de Oliveira e D. Hotza, \textit{Tecnologia de fabricação de revestimentos cerâmicos}, 2. ed. rev. Florianópolis: Editora da UFSC, 2015. ISBN 978-85-328-0710-6.

\bibitem{mctic2017} Ministério da Ciência, Tecnologia, Inovações e Comunicações, \textit{Modelagem setorial de opções de baixo carbono para o setor de cerâmica}. Brasília: MCTIC, 2017.

\bibitem{monfort2010} E. Monfort \textit{et al.}, ``Analysis of energy consumption and carbon dioxide emissions in ceramic tile manufacture'', \textit{Boletín de la Sociedad Española de Cerámica y Vidrio}, vol. 49, n. 4, pp. 303--310, 2010.

\bibitem{alves2023a} C. L. Alves \textit{et al.}, ``Application of flowsheet simulation methodology to improve productivity and sustainability of porcelain tile manufacturing'', \textit{Machines}, vol. 11, n. 2, art. 137, 2023. DOI: 10.3390/machines11020137.

\bibitem{portobello2025} Portobello Grupo, \textit{Relatório de sustentabilidade 2025}. Tijucas: Portobello, 2026. Disponível em: \url{https://www.portobello.com.br/produtos/download/6850/RS2025Portobello_V12.pdf}. Acesso em: 24 ago. 2026.

\bibitem{rigobelo2026} R. Rigobelo, ``Dúvidas técnicas processo produtivo'', mensagem eletrônica a E. A. B. Vaz, 21 ago. 2026. Dados de processo da unidade de Tijucas cedidos pela Portobello para uso acadêmico (comunicação pessoal).

\bibitem{negre2000} F. Negre, J. C. Jarque, C. Feliu e J. E. Enrique, ``Estudo da operação de secagem por atomização de pós cerâmicos em escala industrial: seu controle e automatização'', \textit{Cerâmica Industrial}, vol. 5, n. 2, 2000.

\bibitem{brown1996} H. L. Brown, B. B. Hamel e B. A. Hedman, \textit{Energy analysis of 108 industrial processes}. Lilburn: Fairmont Press, 1996.

\bibitem{cengel2013} Y. A. Çengel e M. A. Boles, \textit{Termodinâmica}, 7. ed. Porto Alegre: AMGH, 2013.

\bibitem{ipcc2006} IPCC, \textit{2006 IPCC guidelines for national greenhouse gas inventories}, vol. 2: Energy. Hayama: IGES, 2006.

\bibitem{mcti2024} Ministério da Ciência, Tecnologia e Inovação, \textit{Fatores de emissão de CO\textsubscript{2} pela geração de energia elétrica no Sistema Interligado Nacional}. Brasília: MCTI, 2024.

\bibitem{della2012} L. P. Della, A. R. Melo e C. R. Melo, ``Análise do sistema de resfriamento final de um forno a rolo do setor cerâmico'', \textit{Cerâmica Industrial}, vol. 17, n. 4, pp. 22--32, 2012.

\bibitem{boyce2012} M. P. Boyce, \textit{Gas turbine engineering handbook}, 4. ed. Oxford: Butter\-worth-Heine\-mann, 2012.

\bibitem{netl2022} National Energy Technology Laboratory, \textit{A literature review of hydrogen and natural gas turbines: current state of the art with regard to performance and NO\textsubscript{x} control}. Pittsburgh: NETL, 2022.

\bibitem{iea2025a} International Energy Agency, \textit{Global Hydrogen Review 2025}. Paris: IEA, 2025.

\bibitem{iea2025b} International Energy Agency, \textit{Global Hydrogen Review 2025: assumptions annex}. Paris: IEA, 2025.

\bibitem{canalsolar2022} Canal Solar, ``Pesquisa da UFSC mapeia produtividade solar em regiões do Estado'', 2022. Disponível em: \url{https://canalsolar.com.br/pesquisa-da-ufsc-mapeia-produtividade-solar-em-regioes-do-estado/}. Acesso em: 1 out. 2026.

\bibitem{pvmag2021} pv magazine USA, ``East-west orientation? Here are two new approaches to designing your solar project'', 9 ago. 2021. Disponível em: \url{https://pv-magazine-usa.com/2021/08/09/east-west-orientation-here-are-two-new-approaches-to-designing-your-solar-project/}. Acesso em: 1 out. 2026.

\bibitem{bolinger2022} M. Bolinger e G. Bolinger, ``Land requirements for utility-scale PV: an empirical update on power and energy density'', \textit{IEEE Journal of Photovoltaics}, vol. 12, n. 2, pp. 589--594, 2022.

\bibitem{aneel2026} Agência Nacional de Energia Elétrica, Resolução Homologatória nº 3.602, de 18 de agosto de 2026: Revisão Tarifária Periódica de 2026 da Celesc Distribuição S.A., \textit{Diário Oficial da União}, seção 1, p. 55, 21 ago. 2026.

\bibitem{aneel2021} Agência Nacional de Energia Elétrica, Resolução Normativa nº 1.000, de 7 de dezembro de 2021: regras de prestação do serviço público de distribuição de energia elétrica.

\bibitem{schlomer2014} S. Schlömer \textit{et al.}, ``Annex III: Technology-specific cost and performance parameters'', em IPCC, \textit{Climate change 2014: mitigation of climate change}. Cambridge: Cambridge University Press, 2014.

\bibitem{lei14948} Brasil, Lei nº 14.948, de 2 de agosto de 2024: institui o marco legal do hidrogênio de baixa emissão de carbono, \textit{Diário Oficial da União}, 2024.

\bibitem{cepci2026} Chemical Engineering, ``Inflationary pressures push 2025 CEPCI annual average higher than 2024'', abr. 2026. Disponível em: \url{https://www.chemengonline.com/inflationary-pressures-push-2025-cepci-annual-average-higher-than-2024/}. Acesso em: 17 nov. 2026.

\bibitem{toweringskills} Towering Skills, ``Cost indices'' (CEPCI anual de 2022 e de junho de 2024, da Chemical Engineering). Disponível em: \url{https://toweringskills.com/financial-analysis/cost-indices/}. Acesso em: 17 nov. 2026.

\bibitem{agrolink2026} Agrolink, ``Dólar fecha em queda e acumula desvalorização no ano'', 27 out. 2026. Disponível em: \url{https://www.agrolink.com.br/noticias/dolar-fecha-em-queda-e-acumula-desvalorizacao-no-ano_520298.html}. Acesso em: 17 nov. 2026.

\bibitem{izidoro2016} D. L. Izidoro, ``Recuperação de calor residual de baixa qualidade em usinas termelétricas de ciclo combinado'', Dissertação (Mestrado em Conversão de Energia), Universidade Federal de Itajubá, Itajubá, 2016, cap. 5.

\bibitem{towler2022} G. Towler e R. Sinnott, \textit{Chemical engineering design}, 3. ed. Elsevier, 2022.

\bibitem{epe2025} Empresa de Pesquisa Energética, \textit{Plano Decenal de Expansão de Energia 2035: parâmetros de custos -- geração e transmissão}. Rio de Janeiro: EPE, 2025.

\bibitem{irena2025} International Renewable Energy Agency, \textit{Renewable power generation costs in 2024}. Abu Dhabi: IRENA, 2025.

\bibitem{nrel2014} National Renewable Energy Laboratory, \textit{Hydrogen station compression, storage, and dispensing: technical status and costs}, NREL/BK-6A10-58564, 2014.

\bibitem{gtw2018} Gas Turbine World, \textit{2018 GTW Handbook}, vol. 33. Pequot Publishing, 2018.

\bibitem{solar2026} Solar Turbines, ``Taurus 60'', página do conjunto gerador (dados nas condições ISO: 5.670~kWe, 11.430~kJ/kWh, escape de 78.385~kg/h a 510~°C). Disponível em: \url{https://www.solarturbines.com/en_US/products/power-generation-packages/taurus-60.html}. Acesso em: 26 nov. 2026.

\bibitem{arup2022} Arup, \textit{Study to develop cost and stock assumptions for hydrogen-ready industrial boilers}, relatório para o Department for Business, Energy and Industrial Strategy (Reino Unido), 2022.

\bibitem{turton2018} R. Turton \textit{et al.}, \textit{Analysis, synthesis, and design of chemical processes}, 5. ed. Prentice Hall, 2018, cap. 7.

\bibitem{venturini2021} O. J. Venturini, ``EEN707 Projeto Integrado I: caderno de encargos e estimativa de custos'', notas de aula, Universidade Federal de Itajubá, 2021.

\bibitem{seider2017} W. D. Seider \textit{et al.}, \textit{Product and process design principles}, 4. ed. Wiley, 2017.

\bibitem{lora2004} E. E. S. Lora e M. A. R. Nascimento, \textit{Geração termelétrica: planejamento, projeto e operação}. Rio de Janeiro: Interciência, 2004, cap. 19.

\bibitem{scgas2026} SCGÁS, ``Entenda a tarifa TG1'' (segmento industrial, cobrança em cascata por faixa de consumo diário), vigente desde 1 jul. 2026 (condições gerais: poder calorífico de referência de 9.400 kcal/m³, 1,033 kgf/cm² e 20 °C; também no \textit{Diário Oficial do Estado de Santa Catarina}, 1 jul. 2026). Disponível em: \url{https://www.scgas.com.br/scgas/site/tarifas-e-simuladores/entenda-a-tarifa-tg1}. Acesso em: 26 nov. 2026.

\bibitem{scgas2026b} SCGÁS, ``Tabela de tarifas para o segmento de matéria-prima'', vigente desde 1 jul. 2026. Disponível em: \url{https://www.scgas.com.br/scgas/site/tarifas-e-simuladores/materia-prima}. Acesso em: 24 nov. 2026.

\bibitem{scgas2026c} SCGÁS, ``Geração e cogeração de energia'' (tarifa TGG, faturamento por faixa), vigente desde 1 jul. 2026. Disponível em: \url{https://www.scgas.com.br/scgas/site/para-a-industria/geracao-e-cogeracao-de-energia}. Acesso em: 26 nov. 2026.

\bibitem{ccee2025} Câmara de Comercialização de Energia Elétrica, ``Preço de Liquidação das Diferenças horário, submercado Sul, 2024 e 2025'', dados abertos (PLD\_HORARIO). Disponível em: \url{https://dadosabertos.ccee.org.br/dataset/pld_horario}. Acesso em: 14 nov. 2026.

\bibitem{pvgis2026} European Commission, Joint Research Centre, \textit{Photovoltaic Geographical Information System (PVGIS)}, versão 5.3, série horária. Disponível em: \url{https://re.jrc.ec.europa.eu/pvg_tools/en/}. Acesso em: 14 nov. 2026.

\bibitem{greener2026} L. Neves, ``Sistemas fotovoltaicos ficam até 7\% mais caros no 1º semestre de 2026, aponta Greener'', \textit{pv magazine Brasil}, 25 set. 2026. Disponível em: \url{http://www.pv-magazine-brasil.com/2026/09/25/sistemas-fotovoltaicos-ficam-ate-7-mais-caros-no-1o-semestre-de-2026-aponta-greener/}. Acesso em: 17 nov. 2026.

\bibitem{yermolenko2026} H. Yermolenko, ``European carbon prices exceeded €90/t in November'', \textit{GMK Center}, nov. 2026. Disponível em: \url{https://gmk.center/en/news/european-carbon-prices-exceeded-e90-t-in-november/}. Acesso em: 17 nov. 2026.

\bibitem{lei15042} Brasil, Lei nº 15.042, de 11 de dezembro de 2024: institui o Sistema Brasileiro de Comércio de Emissões de Gases de Efeito Estufa (SBCE), \textit{Diário Oficial da União}, 2024.

\bibitem{chemeng2026} Chemical Engineering, ``Economic indicators: Chemical Engineering Plant Cost Index'' (tabela anual; CEPCI de 2007). Disponível em: \url{https://www.chemengonline.com/cepci}. Acesso em: 17 nov. 2026.

\bibitem{wartsila} Wärtsilä, ``Technology comparison: derating due to ambient temperature''. Disponível em: \url{https://www.wartsila.com/energy/learn-more/technology-comparison-engines-vs-aeros/derating-due-to-ambient-temperature}. Acesso em: 17 nov. 2026.

\bibitem{embrapa} Embrapa Clima Temperado, Laboratório de Agrometeorologia, ``Temperatura média mensal, Santa Catarina'' (normal 1961--1990, INMET). Disponível em: \url{https://agromet.cpact.embrapa.br/tab/temp_media_SC.pdf}. Acesso em: 17 nov. 2026.

\bibitem{james2024} B. D. James \textit{et al.}, \textit{Hydrogen production cost with anion exchange membrane electrolysis}. Strategic Analysis Inc. para o Department of Energy (EUA), set. 2024. Disponível em: \url{https://www.osti.gov/servlets/purl/2476013}. Acesso em: 25 nov. 2026.

\bibitem{valutafx} ValutaFX, ``Histórico GBP/USD de 2022'' (média anual). Disponível em: \url{https://www.valutafx.com/history/gbp-usd-2022-12-28}. Acesso em: 25 nov. 2026.

\bibitem{anp2026} Agência Nacional do Petróleo, Gás Natural e Biocombustíveis, \textit{Anuário estatístico 2026: fatores de conversão, densidades e poderes caloríficos inferiores} (valores médios de 2025; fonte: EPE, Balanço Energético Nacional). Disponível em: \url{https://www.gov.br/anp/pt-br/centrais-de-conteudo/publicacoes/anuario-estatistico/arquivos-anuario-estatistico-2026/outras-pecas-documentais/fatores-de-conversao-2026.pdf}. Acesso em: 24 nov. 2026.

\bibitem{alves2023b} C. L. Alves \textit{et al.}, ``Improving the sustainability of porcelain tile manufacture by flowsheet simulation'', \textit{Ceramics International}, vol. 49, n. 15, pp. 24581--24597, 2023.

\bibitem{sartor2021} R. P. Sartor, J. A. Speck, A. E. Savi, J. H. Piva e E. G. P. Antunes, ``Estudo da influência da aplicação de pó residual do processo de retífica à seco de placas cerâmicas em argamassas colantes industrializadas'', \textit{Mix Sustentável}, vol. 7, n. 3, pp. 29--40, 2021. DOI: 10.29183/2447-3073.MIX2021.v7.n3.29-40.

\bibitem{castineiras2025} S. L. P. Castiñeiras-Filho, G. Fortunato, S. Cardoso, L. F. M. Frutuoso, F. Pradelle, E. Fagundes de Almeida e E. Fernández y Fernández, ``Techno-economic assessment of the use of green hydrogen: case study in the ceramic industry'', \textit{ENERLAC: Revista de Energía de Latinoamérica y el Caribe}, vol. 8, n. 2, 2024 (publicado em 2025). Disponível em: \url{https://enerlac.olacde.org/index.php/ENERLAC/es/article/view/311}. Acesso em: 26 nov. 2026.

\bibitem{ferrer2019} S. Ferrer, A. Mezquita, V. M. Aguilella e E. Monfort, ``Beyond the energy balance: exergy analysis of an industrial roller kiln firing porcelain tiles'', \textit{Applied Thermal Engineering}, vol. 150, pp. 1002--1015, 2019. DOI: 10.1016/j.applthermaleng.2019.01.052.

\bibitem{wang2020} Y. Wang, H. Yang e K. Xu, ``Thermal performance combined with cooling system parameters study for a roller kiln based on energy-exergy analysis'', \textit{Energies}, vol. 13, n. 15, 3922, 2020. DOI: 10.3390/en13153922.

\bibitem{bejan1996} A. Bejan, G. Tsatsaronis e M. Moran, \textit{Thermal design and optimization}. New York: Wiley, 1996.

\bibitem{lazzaretto2006} A. Lazzaretto e G. Tsatsaronis, ``SPECO: a systematic and general methodology for calculating efficiencies and costs in thermal systems'', \textit{Energy}, vol. 31, n. 8--9, pp. 1257--1289, 2006.

\bibitem{mezquita2012} A. Mezquita, E. Monfort, E. Vaquer, S. Ferrer, M. A. Arnal, J. Toledo e M. A. Cuesta, ``Optimización energética en la fabricación de baldosas cerámicas mediante el uso de aceite térmico'', \textit{Boletín de la Sociedad Española de Cerámica y Vidrio}, vol. 51, n. 4, pp. 183--190, 2012. DOI: 10.3989/cyv.272012.

